% ============================================================
\subsubsection{Arquitectura C4 y configuración del RAG}
% ============================================================

La arquitectura se documentó con el modelo C4~\cite{brown2024c4}, que describe
el sistema en cuatro niveles de abstracción.

\paragraph{Nivel 1: Contexto}

NutriAPP es el sistema central (Figura~\ref{fig:c4_contexto}). El nutricionista
accede desde un navegador web y la única dependencia externa es la API de Groq,
que provee la inferencia del modelo de lenguaje (RNF05).

\figura{c4_contexto.png}{Diagrama C4, nivel 1: contexto del sistema}{fig:c4_contexto}

\paragraph{Nivel 2: Contenedores}

El sistema se compone de tres contenedores (Figura~\ref{fig:c4_contenedores}):
una aplicación de página única en React servida por Vite, una API REST en
Express con TypeScript y una base de datos PostgreSQL accedida mediante Prisma ORM.

\figura{c4_contenedores.png}{Diagrama C4, nivel 2: contenedores}{fig:c4_contenedores}

\paragraph{Nivel 3: Componentes del backend}

El backend se organiza en cuatro routers con su controlador
(Tabla~\ref{tab:c4_comp} y Figura~\ref{fig:c4_componentes}). El servicio
\texttt{aiService} concentra la comunicación con la API de Groq y el middleware
de autenticación verifica el token JWT en cada ruta protegida (RNF02).

\begin{table}[H]
\centering\small
\caption{Componentes del backend (nivel 3 C4)}\label{tab:c4_comp}
\begin{tabularx}{\linewidth}{|l|l|X|}
\hline
\textbf{Router} & \textbf{Controlador} & \textbf{Responsabilidad} \\ \hline
\texttt{/api/auth} & \texttt{authController} & Inicio de sesión y administración de nutricionistas \\ \hline
\texttt{/api/patients} & \texttt{patientController} & Pacientes, historia clínica, biometría, antropometría, hábitos dietéticos y menú semanal \\ \hline
\texttt{/api/foods} & \texttt{foodController} & Catálogo de alimentos y grupos SMAE \\ \hline
\texttt{/api/ai} & \texttt{aiController} & Chat, recomendaciones, análisis de alimentos, planes de comida, carga de documentos e historial \\ \hline
\end{tabularx}
\end{table}

\figura{c4_componentes.png}{Diagrama C4, nivel 3: componentes del backend}{fig:c4_componentes}

\paragraph{Nivel 4: Código}

La Figura~\ref{fig:c4_codigo} detalla el servicio \texttt{aiService}, que expone
los métodos \texttt{chat}, \texttt{chatWithDocument},
\texttt{generateNutritionRecommendation}, \texttt{generateMealPlan} y
\texttt{analyzeFood}, además de la función interna \texttt{buildSmaeContext}.

\figura{c4_codigo.png}{Diagrama C4, nivel 4: código del servicio de IA}{fig:c4_codigo}

\paragraph{Configuración del pipeline RAG}

El asistente combina tres fuentes de contexto antes de cada llamada al modelo:

\begin{enumerate}[leftmargin=2em]
  \item \textbf{Datos clínicos del paciente}: historial, antropometría,
        biometría y hábitos del paciente seleccionado (RF09) o de todos los
        pacientes en el chat global (RF10).
  \item \textbf{Catálogo SMAE}: cuando la consulta contiene términos como
        \emph{plan}, \emph{alimento}, \emph{menú}, \emph{equivalente} o
        \emph{kcal}, el sistema recupera de la base de datos ocho alimentos por
        cada uno de siete grupos de equivalentes (56 alimentos) y los añade al
        prompt con su porción, energía, macronutrientes e índice glucémico.
  \item \textbf{Documento del nutricionista}: el texto extraído del PDF cargado
        (RF11), limitado a los primeros 6000 caracteres.
\end{enumerate}

\pendiente{precisa en el texto qué entiendes por «recuperación»: hoy el sistema
selecciona alimentos por grupo y trunca el PDF, no hace búsqueda semántica con
embeddings. Descríbelo tal cual o implementa la búsqueda antes de redactarlo}.

La Tabla~\ref{tab:rag_param} resume los parámetros configurados.

\begin{table}[H]
\centering\small
\caption{Parámetros de configuración del asistente de IA}\label{tab:rag_param}
\begin{tabularx}{\linewidth}{|p{5cm}|X|}
\hline
\textbf{Parámetro} & \textbf{Valor} \\ \hline
Proveedor de inferencia & API de Groq (compatible con OpenAI) \\ \hline
Modelo por defecto & \pendiente{confirma: el código usa \texttt{openai/gpt-oss-120b}; la metodología dice Llama~3.3~70B} \\ \hline
Modelos seleccionables & \texttt{gpt-oss-120b}, \texttt{llama-3.3-70b-versatile}, \texttt{llama-3.1-8b-instant} \\ \hline
Temperatura & 0.3 \\ \hline
Máximo de tokens de respuesta & 8000 \\ \hline
Historial enviado al modelo & 20 mensajes de la conversación \\ \hline
Contexto SMAE & 7 grupos $\times$ 8 alimentos, activado por palabras clave \\ \hline
Documentos admitidos & PDF o texto plano, hasta 10\,MB \\ \hline
Texto del documento usado & primeros 6000 caracteres \\ \hline
Extracción de texto & \texttt{pdf-parse} 2.x \\ \hline
\end{tabularx}
\end{table}

El comportamiento del asistente se restringe mediante un prompt del sistema
(\texttt{backend/prompts/system.md}) que define su alcance clínico, rechaza
solicitudes fuera del dominio nutricional, obliga a presentar los valores en
tablas y declara las fórmulas y referencias admitidas (Mifflin-St~Jeor,
Harris-Benedict, FAO/OMS, SMAE 4.ª ed.).

\begin{lstlisting}[caption={Construcción del contexto SMAE por grupo (extracto de \texttt{aiService.ts})},label={lst:smae}]
async function buildSmaeContext(grupos?: string[]): Promise<string> {
  const targetGroups = grupos ?? [
    'Verduras', 'Frutas', 'Cereales', 'Leguminosas', 'AOA', 'Leche', 'Aceites y grasas',
  ];
  const perGroup = 8;
  const perGroupResults = await Promise.all(
    targetGroups.map(g =>
      prisma.food.findMany({
        where: { source: 'SMAE', grupoSmae: g },
        select, orderBy: { name: 'asc' }, take: perGroup,
      })
    )
  );
  // ... se formatea cada alimento como:
  // "- Nombre (porción): kcal P:g G:g HC:g Fib:g IG"
}
\end{lstlisting}
